\documentclass[12pt,a4paper]{article}
        \makeatletter
        \def\input@path{{../}}
        \makeatother
        \usepackage{almeria-fichas}

        \FichaMeta{Sesion 15}{Comparar graficas}{Bloque 3 · Datos, funciones y graficas}{Comparar dos graficas para descubrir relaciones entre datos y decidir cual explica mejor una situacion turistica.}{Conceptual}{Analizar, Evaluar, Crear}

        \begin{document}

        \FichaCabecera

        \begin{materialesbox}
        \begin{itemize}
    \item Ficha impresa
    \item Lapiz
    \item Regla
    \item Calculadora si se necesita
\end{itemize}
        \end{materialesbox}

        \begin{pasosbox}
        \begin{enumerate}
    \item Observa primero cada grafica por separado.
    \item Anota maximos, minimos y tendencia general.
    \item Despues compara lo que ocurre a la vez en las dos graficas.
    \item Escribe una conclusion bien justificada.
\end{enumerate}
        \end{pasosbox}

        \begin{recordatoriobox}
        \begin{itemize}
    \item Comparar no es solo mirar el dibujo: hay que usar datos.
    \item Dos variables pueden cambiar a la vez, pero eso no siempre significa causa.
    \item Una buena comparacion cita semejanzas y diferencias.
\end{itemize}
        \end{recordatoriobox}

        \begin{apoyobox}
        \begin{itemize}
    \item Puedes hacer una tabla con dos columnas: grafica A y grafica B.
    \item Subraya los meses o anos que aparecen en ambas.
    \item Escribe primero los datos y despues tu opinion.
\end{itemize}
        \end{apoyobox}

        \BloomSection{analizar}

\BloomExercise{analizar}{1}{Semejanzas y diferencias}{Compara una grafica de lluvia y otra de aforo de playa. Escribe dos semejanzas y dos diferencias.}{7}

\BloomExercise{analizar}{2}{Relacion posible}{Analiza si cuando sube una variable baja la otra y que tipo de relacion parece haber.}{7}

\BloomSection{evaluar}

\BloomExercise{evaluar}{1}{Grafica mas clara}{Evalua cual de dos graficas explica mejor una ruta turistica y justifica tu respuesta con dos criterios.}{7}

\BloomExercise{evaluar}{2}{Escala enganosa}{Evalua si una grafica con el eje vertical recortado ayuda o puede confundir al lector.}{7}

\BloomSection{crear}

\BloomExercise{crear}{1}{Mejora de presentacion}{Diseña una version mejorada de una de las dos graficas indicando titulo, ejes y escala.}{8}

\BloomExercise{crear}{2}{Comentario comparativo}{Redacta un pequeno comentario final de 5 o 6 lineas para la guia en el que compares dos graficas de la ciudad.}{8}

        \begin{evidenciabox}
        \begin{itemize}
    \item Comparar dos representaciones con datos concretos.
    \item Argumentar si una grafica es mas util que otra.
    \item Proponer una mejora de presentacion.
\end{itemize}
        \end{evidenciabox}

        \end{document}
